\documentclass[10pt,a4paper]{amsart}
\usepackage[margin=19mm]{geometry}
\usepackage{amsmath,amssymb,mathtools}
\usepackage[scr=rsfs,cal=euler]{mathalfa}
\usepackage{xfrac}
\usepackage[unicode,hidelinks,bookmarks=false]{hyperref}
\newcommand{\ie}{i.e.\ }
\newcommand{\cf}{cf.\ }
\newcommand{\resp}{respectively\ }
\newcommand{\Subsection}{\subsection}
\providecommand{\nobreakdash}{\nobreak\hbox{-}}
\setlength{\emergencystretch}{2em}
\multlinegap0pt
\usepackage{bm}
\usepackage[shortlabels]{enumitem}
\usepackage{amssymb}
\usepackage{mathtools}
\usepackage{xcolor}
\usepackage[normalem]{ulem}
\newtheorem{remark}{Remark}
\newtheorem{lemma}{Lemma}
\usepackage{cleveref}
\crefname{remark}{Remark}{Remarks}
\crefname{lemma}{Lemma}{Lemmas}
\newcommand{\Div}{\text{div}}
\newcommand{\E}{\ensuremath{\mathbb{E}}}
\newcommand{\R}{\mathbb{R}}
\newcommand{\T}{\mathbb{T}}
\title{Turbulent stretching of FENE dumbbell polymer model via special stochastic scaling and singular limits}
\author{Federico Butori, Yassine Tahraoui}
\date{Source: arXiv:2605.15742v1 (CC BY 4.0)}
\begin{document}
\maketitle

\noindent\emph{Source and licence.} Turbulent stretching of FENE dumbbell polymer model via special stochastic scaling and singular limits. Version arXiv:2605.15742v1 (CC BY 4.0), distributed by arXiv under the Creative Commons Attribution 4.0 International licence, http://creativecommons.org/licenses/by/4.0/, as recorded in the arXiv licence metadata for this exact version and shown on the abstract page https://arxiv.org/abs/2605.15742v1. Full licence notice: \texttt{licenses/arxiv-2605.15742-CC-BY-4.0.txt}.\par\smallskip

\noindent\textit{Editorial scope.} Counted unit: the article's lemma sec:estimates (source0.tex lines 740--746). The article's complete original proof of it is delivered with it (source0.tex lines 747--816, one \texttt{proof} environment of 70 lines). No mathematics is written, repaired or re-derived: the statement and the proof are the article's own lines, in the article's own notation, and no retained line is substituted or re-flowed.  Retained uncounted as the source-local context the proof needs: lines 426--431; lines 454--457; lines 732--736.  Nothing was omitted from any retained span, so the package carries no omission record.  No retained line contains a citation, so the wrapper carries no bibliography and no bibliography entry.  No retained line contains a figure environment, an image-inclusion command or drawing code, so no image is delivered and the package declares no asset dependency.  Editorial formatting only, in addition: an AMS \texttt{amsart} preamble that restates the article's own declarations and macros used by the retained lines, namely the environments \texttt{remark}, \texttt{lemma} on separate counters, together with the standard packages those lines need.  The retained environments print in this document as Remark 1, Lemma 1 in order of appearance, whereas the article prints the same items with the numbering of its own full document.  The complete native source of the article as distributed in the arXiv e-print for this exact version is bundled unchanged as \texttt{sources/200676/source0.tex}.
\begin{align}\label{Def-sigma-k}
\sigma_{k}^{N,\tau}\left(  x\right)     =\theta_{k}^{N,\tau}\frac{k^{\perp}}{\left\vert
k\right\vert }\cos k\cdot x,\qquad k\in K_{+},  \quad 
\sigma_{k}^{N,\tau}\left(  x\right)     =\theta_{k}^{N,\tau}\frac{k^{\perp}}{\left\vert
k\right\vert }\sin k\cdot x,\qquad k\in K_{-}%
\end{align}

      \begin{remark}\label{rmq-bound_A^N}
            For every $N>1$ due to the choice of the coefficients $\sigma_k^N$ (see \eqref{Def-sigma-k}), there exists a constant $C_A$ independent of $N, r\in B(0, 1)$ and $x\in \T^2$ such that for every $v\in \R^2$ it holds
            $$|A^N(r)v| \le C_A|v|.$$
              \end{remark}

\begin{align}\label{approx-linear*SDE}
    &(f_m(t),w_{i,l})_H-(f_{0},w_{i,l})_H=\int_0^t(\mathcal{Y}^mf_m,w_{i,l})_Hds, \quad 1\leq i,l \leq m \\
    &	-\int_0^t\dfrac{1}{\beta}\big(\mathbf{M}\nabla_r(\dfrac{f_m}{\mathbf{M}}),\nabla_r(\dfrac{w_{i,l}}{\mathbf{M}})\big)ds -\alpha_N\int_0^t(\nabla_x f_m,\nabla_x w_{i,l})_Hds- \dfrac{1}{2}\int_0^t\big(\varphi_\varepsilon^2A^N(r) \nabla_rf_m,\nabla_r(\dfrac{w_{i,l}}{\mathbf{M}})\big)ds\notag\\
            &-\sum_{k\in K}\int_0^t(\sigma_k^N.\nabla_xf_mdW^k,w_{i,l})_H -\sum_{k\in K}\int_0^t(\varphi_\varepsilon(\nabla \sigma_k^Nr).\nabla_rf_m dW^k,w_{i,l})_H, \notag
\end{align}

\begin{lemma}\label{sec:estimates}
The solution $f_m$ of \eqref{approx-linear*SDE} satisfies
    \begin{align}\label{key-est-Galerkin}
 \E\Big[\sup_{v\in[0,t]} \Vert f_m(v)\Vert_H^2  \Big]&+\dfrac{4}{\beta}\E\int_0^t [f_m(s)]_V^2ds\leq 2\Vert f_{0}\Vert_H^2\exp{(4\mathcal{K}_\varepsilon t)} \text{ for all } t\in [0,T].
\end{align}
where $\mathcal{K}_\varepsilon=\dfrac{2C_B^2\mathbf{C}(\kappa+1)^2 a_\tau^2+\kappa (5+\kappa)C_A}{4\varepsilon^2}.$
\end{lemma}

\begin{proof} 
By applying It\^o formula and \eqref{approx-linear*SDE}, one gets
  \begin{align*}
   \dfrac{1}{2} \Vert f_m(t)\Vert_H^2&-\dfrac{1}{2}\Vert P_mf_{0}\Vert_H^2=\int_0^t(\mathcal{Y}^mf_m,f_m)_Hds  \notag\\
    &	-\int_0^t\dfrac{1}{\beta}\big(\mathbf{M}\nabla_r(\dfrac{f_m}{\mathbf{M}}),\nabla_r(\dfrac{f_m}{\mathbf{M}})\big)ds -\alpha_N\int_0^t(\nabla_x f_m,\nabla_x f_m)_Hds \notag\\&- \dfrac{1}{2}\int_0^t\big(\sum_{k\in K}\left(\varphi_\varepsilon^2(\nabla \sigma_k^Nr) \otimes (\nabla \sigma_k^Nr)\right) \nabla_rf_m,\nabla_r(\dfrac{f_m}{\mathbf{M}})\big)ds\\
            &-\sum_{k\in K}\int_0^t(\sigma_k^N.\nabla_xf_mdW^k,f_m)_H -\sum_{k\in K}\int_0^t(\varphi_\varepsilon(\nabla \sigma_k^Nr).\nabla_rf_m dW^k,f_m)_H \\
            &+\dfrac{1}{2}[\sum_{i,l=1}^m\sum_{k\in K}\int_0^t(\sigma_k^N.\nabla_xf_m,w_{i,l})_H^2 ds +\int_0^t(\varphi_\varepsilon(\nabla \sigma_k^Nr).\nabla_rf_m,w_{i,l})_H^2 ds\\
            &+2\sum_{k\in K}\int_0^t\big(\sigma_k^N.\nabla_xf_m, \sum_{i,l=1}^m(\varphi_\varepsilon(\nabla \sigma_k^Nr).\nabla_rf_m,w_{i,l})_Hw_{i,l}\big)_H ds].
\end{align*}
Note that 
\begin{align*}
    &\sum_{k\in K}\int_0^t\big(\sigma_k^N.\nabla_xf_m, \sum_{i,l=1}^m(\varphi_\varepsilon(\nabla \sigma_k^Nr).\nabla_rf_m,w_{i,l})_Hw_{i,l}\big)_H ds\\&=\sum_{k\in K}\int_0^t\big(\sigma_k^N.\nabla_xf_m, P_m(\varphi_\varepsilon(\nabla \sigma_k^Nr).\nabla_rf_m)\big)_H ds=-\int_0^t(\mathcal{Y}^mf_m,f_m)_Hds.
\end{align*}
Moreover, we have 
\begin{align*}
    &\dfrac{1}{2}\sum_{i,l=1}^m\sum_{k\in K}\int_0^t(\sigma_k^N.\nabla_xf_m,w_{i,l})_H^2 ds-\alpha_N\int_0^t(\nabla_x f_m,\nabla_x f_m)_Hds\\&=\dfrac{1}{2}\sum_{k\in K}\int_0^t\Vert P_m\sigma_k^N.\nabla_xf_m\Vert_H^2 ds-\alpha_N\int_0^t(\nabla_x f_m,\nabla_x f_m)_Hds\\
    &\leq  \dfrac{1}{2}\sum_{k\in K}\int_0^t\Vert \sigma_k^N.\nabla_xf_m\Vert_H^2 ds-\alpha_N\int_0^t(\nabla_x f_m,\nabla_x f_m)_Hds\leq 0.
\end{align*}
Since $\Div \sigma_k^N=0,$ we get $(\sigma_k^N.\nabla_xf_m,f_m)_H=0$ P-a.s. Therefore
\begin{align}
   \dfrac{1}{2} \Vert f_m(t)\Vert_H^2&+\dfrac{1}{\beta}\int_0^t\big(\mathbf{M}\nabla_r(\dfrac{f_m}{\mathbf{M}}),\nabla_r(\dfrac{f_m}{\mathbf{M}})\big)ds-\dfrac{1}{2}\Vert P_mf_{0}\Vert_H^2\le\notag\\
  &- \dfrac{1}{2}\int_0^t\big(\sum_{k\in K}\left(\varphi_\varepsilon^2(\nabla \sigma_k^Nr) \otimes (\nabla \sigma_k^Nr)\right) \nabla_rf_m,\nabla_r(\dfrac{f_m}{\mathbf{M}})\big)ds\label{ineq-estm-Galer}\\
            & -\sum_{k\in K}\int_0^t(\varphi_\varepsilon(\nabla \sigma_k^Nr).\nabla_rf_m dW^k(s),f_m)_H+ \sum_{k\in K}\sum_{j,k \le m}\dfrac{1}{2}\int_0^t(\varphi_\varepsilon(\nabla \sigma_k^Nr).\nabla_rf_m,w_{i,l})_H^2 ds.\notag
\end{align}
Notice that 
\begin{align}
    &- \dfrac{1}{2}\int_0^t\big(\sum_{k\in K}\left(\varphi_\varepsilon^2(\nabla \sigma_k^Nr) \otimes (\nabla \sigma_k^Nr)\right) \nabla_rf_m,\nabla_r(\dfrac{f_m}{\mathbf{M}})\big)ds+ \sum_{k\in K}\sum_{j,k \le m}\dfrac{1}{2}\int_0^t(\varphi_\varepsilon(\nabla \sigma_k^Nr).\nabla_rf_m,w_{i,l})_H^2 ds \notag\\
         &=- \dfrac{1}{2}\int_0^t\big(\sum_{k\in K}\Vert \varphi_\varepsilon (\nabla \sigma_k^Nr)\cdot\nabla_rf_m\Vert_H^2ds+ \sum_{k\in K}\dfrac{1}{2}\int_0^t\Vert P_m\varphi_\varepsilon(\nabla \sigma_k^Nr).\nabla_rf_m\Vert_H^2 ds\notag\\&-\dfrac{1}{4}\int_0^t\big(\sum_{k\in K}\left(\varphi_\varepsilon^2(\nabla \sigma_k^Nr) \otimes (\nabla \sigma_k^Nr)\right) \nabla_rf_m^2,\nabla(\dfrac{1}{\mathbf{M}})\big)ds\notag\\
        &\leq  -\dfrac{1}{4}\int_0^t\big(\sum_{k\in K}\left(\varphi_\varepsilon^2(\nabla \sigma_k^Nr) \otimes (\nabla \sigma_k^Nr)\right) \nabla_rf_m^2,\nabla(\dfrac{1}{\mathbf{M}})\big)ds\notag\\
        &\leq \dfrac{1}{4}\int_0^t\big(f_m^2, \sum_{k\in K} \Div_r\Big[\left(\varphi_\varepsilon^2(\nabla \sigma_k^Nr) \otimes (\nabla \sigma_k^Nr)\right)\nabla(\dfrac{1}{\mathbf{M}})\Big]\big)ds \leq  \dfrac{\kappa (5+\kappa)C_A}{4\varepsilon^2}\int_0^t\Vert f_m(s)\Vert_H^2 ds. \label{esti-cov-uniq}
\end{align}
Indeed,  by using \Cref{rmq-bound_A^N},  $\varphi_\varepsilon\dfrac{1}{1-\vert r\vert^2} \leq \dfrac{1}{\varepsilon}$  and  
\begin{align*} 
    &\int_0^t\big(f_m^2, \sum_{k\in K} \Div_r\Big[\left(\varphi_\varepsilon^2(\nabla \sigma_k^Nr) \otimes (\nabla \sigma_k^Nr)\right)\nabla(\dfrac{1}{\mathbf{M}})\Big]\big)ds\\&=2\int_0^t\big(f_m^2, \sum_{k\in K} \left(\varphi_\varepsilon \nabla \varphi_\varepsilon(\nabla \sigma_k^Nr) \otimes (\nabla \sigma_k^Nr)\right)\nabla(\dfrac{1}{\mathbf{M}})\big)ds\\
    &+ \int_0^t\big(\varphi_\varepsilon^2 f_m^2, \sum_{k\in K} \left((\nabla \sigma_k^Nr) \otimes (\nabla \sigma_k^Nr)\right):\mathcal{D}^2(\dfrac{1}{\mathbf{M}})\big)ds,
\end{align*}
where $\mathcal{D}^2$ denotes the Hessian matrix given by 
\begin{align*}
    \nabla(\dfrac{1}{\mathbf{M}})=\dfrac{\kappa r}{1-\vert r\vert^2}\dfrac{1}{\mathbf{M}} \text{ and }  \mathcal{D}^2(\dfrac{1}{\mathbf{M}})=\dfrac{\kappa(1-\vert r\vert^2)I+\kappa(2+\kappa)r\otimes r}{(1-\vert r\vert^2)^2}\dfrac{1}{\mathbf{M}}.
\end{align*}
Finally, note that
\begin{align*}
   &- \sum_{k\in K}\int_0^t(\varphi_\varepsilon(\nabla \sigma_k^Nr).\nabla_rf_m ,f_m)_H dW^k(s)=\dfrac{1}{2}\sum_{k\in K}\int_0^t(f_m^2 ,(\nabla \sigma_k^Nr)\cdot\nabla(\varphi_\varepsilon \dfrac{1}{\mathbf{M}})) dW^k(s)\\
   &=\dfrac{1}{2}\sum_{k\in K}\int_0^t(\dfrac{f_m^2}{\mathbf{M}} ,(\nabla \sigma_k^Nr)\cdot\nabla\varphi_\varepsilon) dW^k(s)+\dfrac{1}{2}\sum_{k\in K}\int_0^t(\dfrac{f_m^2}{\mathbf{M}}, \varphi_\varepsilon(\nabla \sigma_k^Nr)\cdot \dfrac{\kappa r}{1-\vert r\vert^2} ) dW^k(s).
\end{align*}
We have
\begin{align*}
    & \sum_{k\in K} \vert(\nabla \sigma_k^Nr)\cdot\nabla\varphi_\varepsilon\vert^2 \leq \dfrac{1}{\varepsilon^2} \sum_{k\in K}\vert (\nabla \sigma_k^Nr)\vert^2
     \leq \dfrac{a_\tau^2}{\varepsilon^2} \sum_{k\in K} \dfrac{1}{\vert k\vert^2}\leq  \dfrac{ \mathbf{C}a_\tau^2}{\varepsilon^2}, \\
     & \sum_{k\in K} \vert\varphi_\varepsilon(\nabla \sigma_k^Nr)\cdot \dfrac{\kappa r}{1-\vert r\vert^2}\vert^2 \leq \dfrac{\kappa^2}{\varepsilon^2} \sum_{k\in K}\vert (\nabla \sigma_k^Nr)\vert^2
     \leq \dfrac{\kappa^2 a_\tau^2}{\varepsilon^2} \sum_{k\in K} \dfrac{1}{\vert k\vert^2}\leq  \dfrac{\mathbf{C}\kappa^2 a_\tau^2}{\varepsilon^2}
\end{align*}
since $\sum_{k\in K} \dfrac{1}{\vert k\vert^2}\leq \mathbf{C}$ and $\mathbf{C}$ is independent of $N$. Next, Burkholder-Davis-Gundy inequality ensures the existence of $C_B>0$ such that 
\begin{align*}
  &\E \sup_{v\in[0,t]}   \vert \sum_{k\in K}\int_0^v(\varphi_\varepsilon(\nabla \sigma_k^Nr).\nabla_rf_m ,f_m)_H dW^k(s)\vert \\
  &\leq  \dfrac{C_B}{2}\E\Big[\big(\sum_{k\in K}\int_0^t(\dfrac{f_m^2}{\mathbf{M}} ,(\nabla \sigma_k^Nr)\cdot\nabla\varphi_\varepsilon)^2 ds\big)^{\frac{1}{2}} \Big]+\dfrac{C_B}{2}\E\Big[(\sum_{k\in K}\int_0^t(\dfrac{f_m^2}{\mathbf{M}}, \varphi_\varepsilon(\nabla \sigma_k^Nr)\cdot \dfrac{\kappa r}{1-\vert r\vert^2} )^2 ds\big)^{\frac{1}{2}} \Big]\\
  &\leq  \dfrac{C_B\sqrt{\mathbf{C}}(\kappa+1) a_\tau}{2\varepsilon}\E\Big[(\int_0^t(\dfrac{f_m^2}{\mathbf{M}}, 1)^2 ds\big)^{\frac{1}{2}} \Big] \\ 
  &\leq \dfrac{C_B\sqrt{\mathbf{C}}(\kappa+1) a_\tau}{2\varepsilon}\E\Big[(\int_0^t\Vert f_m(s)\Vert_H^4 ds\big)^{\frac{1}{2}} \Big]\\
  &\leq  \dfrac{C_B\sqrt{\mathbf{C}}(\kappa+1) a_\tau}{2\varepsilon}\E\Big[\sup_{v\in[0,t]} \Vert f_m(v)\Vert_H (\int_0^t\Vert f_m(s)\Vert_H^2 ds\big)^{\frac{1}{2}} \Big].
\end{align*} 
Young's inequality ensures
\begin{align*}
     &\E \sup_{v\in[0,t]}   \vert \sum_{k\in K}\int_0^v(\varphi_\varepsilon(\nabla \sigma_k^Nr).\nabla_rf_m ,f_m)_H dW^k(s)\vert \\&\leq  \dfrac{C_B^2\mathbf{C}(\kappa+1)^2 a_\tau^2}{2\varepsilon^2}\E\Big[ \int_0^t\Vert f_m(s)\Vert_H^2 ds \Big]+ \dfrac{1}{4}\E\Big[\sup_{v\in[0,t]} \Vert f_m(v)\Vert_H^2  \Big].
\end{align*}
Therefore, from \eqref{ineq-estm-Galer}, we infer for any $t\in [0,T]$ \begin{align*}
 \dfrac{1}{4}\E\Big[\sup_{v\in[0,t]} \Vert f_m(v)\Vert_H^2  \Big]&+\dfrac{1}{\beta}\E\int_0^t [f_m(s)]_V^2ds\leq \dfrac{1}{2}\Vert f_{0}\Vert_H^2
 +\mathcal{K}_\varepsilon\E\Big[ \int_0^t\Vert f_m(s)\Vert_H^2 ds \Big]
\end{align*}
where $\mathcal{K}_\varepsilon=\dfrac{2C_B^2\mathbf{C}(\kappa+1)^2 a_\tau^2+\kappa (5+\kappa)C_A}{4\varepsilon^2}.$   Grönwall's lemma then ensures \eqref{key-est-Galerkin}. 
\end{proof}

\end{document}
